 %============================================================
% MTH 112 Project - Front Matter
% Chapter 0 - Front Matter
% Updated 202504
%============================================================

\setcounter{chapter}{-1} 
\chapter{Front Matter}  \label{chapter-frontmatter}

\section*{Colophon}
{\bf Project Lead} Ross Kouzes\\[0.5em]
{\bf Contributing Authors} Shane Horner, Ross Kouzes, Scot Leavitt\\[0.5em]
{\bf Additional Contributors} Amy Cakebread, Bethany Downs, Wendy Fresh, Dave Froemke, Peter Haberman, Shane Horner, Scot Leavitt, Emily Nelson, Kim Neuburger,Emily O’Sullivan,  Jeff Pettit, Dennis Reynolds, Rebecca Ross, Heiko Spoddeck, Greta Swanson, Stephanie Yurasits \\[0.5em]
\copyright 2025 Portland Community College\\[0.5em]
This work is licensed under the Creative Commons Attribution-Noncommercial-ShareAlike (CC BY-NC-SA) 4.0 International License. To view a copy of this license, visit \href{CreativeCommons.org}{https://creativecommons.org/licenses/by-nc-sa/4.0/}\\


\newpage
\vspace{0.75in}

This is the second part of the two part lab manual for MTH 111 and MTH 112 at Portland Community College. The first chapter is \say{Chapter 5} to align with \href{http://tiny.cc/111Z-Supplement}{the first portion of the lab manual}.

\vspace{0.5em}

This lab manual was designed with the intent of being used in the following manner.  In each section:

\setlist{itemsep=6pt}
\begin{itemize}
	\item Students will complete the Preparation Exercises before receiving instruction on the content.  Some instructors will have students do this before coming to class, while others might do it as a warmup in class.
	\item There will be some instructor-led presentation of the content.  This may be a formal class lecture, a discovery activity, a video lecture, or something else.
	\item Students will then engage in Practice Exercises in a group setting to reinforce their initial understanding of the foundational concepts.
	\item An instructor will then assign one or more group activities.  Because many of these activities are web-based and instructors can choose to use different activities, we have not included any of the possible activities in this document. 
	\item The Definitions are meant to provide a single location for all definitions and some key concepts.  Students can use this as a resource after having covered the topics in class.
	\item Students will complete some or all of the Exit Exercises, as decided by their instructor, to summarize their understanding of key concepts.
	\item In general, not every section of this lab manual will take the same amount of class time.  Some sections might be half-day topics, while others could be two-day topics.	
	\item Each instructor will identify for their class which exercises will be submitted as part of the course grade.
\end{itemize}
~

Course Resource Links:\\
{\it Algebra and Trigonometry 2e}:  \href{http://tiny.cc/111Z-Textbook}{http://tiny.cc/112Z-Textbook}\\
MTH 112 Supplement: \href{http://tiny.cc/111Z-Supplement}{http://tiny.cc/112Z-Supplement}

\newpage
~